% Punjabi Shahmukhi adaptation. Original mathematical macro meanings: OpenLogic,
% revision 9620cc73f9c8e0ad003c514a5d3748f29611c4c0, CC BY 4.0.
\documentclass[12pt,a4paper,oneside,openany]{book}
\usepackage[margin=23mm,headheight=16pt]{geometry}
\usepackage{amsmath,amssymb,nicefrac}
\usepackage{tikz,xcolor}
\colorlet{oldiagcolorA}{black}
\colorlet{oldiagcolorB}{gray}
\definecolor{oldiagcolorC}{HTML}{a81c21}
\usepackage{fontspec}
\usepackage[unicode,hidelinks]{hyperref}
\usepackage{polyglossia}
\setdefaultlanguage{english}
\setotherlanguage{arabic}
\setmainfont{Noto Serif}
\newfontfamily\arabicfont[Script=Arabic]{Noto Naskh Arabic}
\newfontfamily\arabicfontsf[Script=Arabic]{Noto Naskh Arabic}
\newfontfamily\arabicfonttt[Script=Arabic]{Noto Naskh Arabic}
\usepackage{setspace,needspace}
\hypersetup{pdftitle={OpenLogic: Sets - Punjabi Shahmukhi},pdfauthor={Open Logic Project; machine translation by Codex},pdfsubject={Provisional Punjabi Shahmukhi translation; seven source units, not the complete edition}}
\setstretch{1.15}
\setlength{\parindent}{0pt}
\setlength{\parskip}{0.45em}
% Arabic scripts do not have a native italic style in these families.
% Preserve emphasis with bold upright rather than silently losing font shape.
\renewcommand{\emph}[1]{{\upshape\bfseries #1}}
\widowpenalty=10000
\clubpenalty=10000
\displaywidowpenalty=10000
\binoppenalty=10000
\relpenalty=10000
\emergencystretch=3em
\pagestyle{plain}
\let\phi\varphi
\newcommand{\Setabs}[2]{\{#1:#2\}}
\newcommand{\Pow}[1]{\wp(#1)}
\newcommand{\Nat}{\mathbb{N}}
\newcommand{\Int}{\mathbb{Z}}
\newcommand{\Rat}{\mathbb{Q}}
\newcommand{\Real}{\mathbb{R}}
\newcommand{\PosInt}{\mathbb{Z}^{+}}
\newcommand{\Bin}{\mathbb{B}}
\newcommand{\len}[1]{\operatorname{len}(#1)}
\newcommand{\tuple}[1]{\langle #1\rangle}
\newcommand{\lif}{\to}
\newenvironment{tagblock}[1]{}{}
\newenvironment{explain}{}{}
\newenvironment{digress}{\par}{\par}
\newcounter{statement}[section]
\renewcommand{\thestatement}{\thesection.\arabic{statement}}
\newcommand{\pnbstatement}[2]{\par\Needspace{4\baselineskip}\refstepcounter{statement}\textbf{#1 \textenglish{\thestatement}۔ #2}\par\nobreak}
\newenvironment{defn}[1][]{\pnbstatement{تعریف}{#1}}{\par}
\newenvironment{ex}[1][]{\pnbstatement{مثال}{#1}}{\par}
\newenvironment{prop}[1][]{\pnbstatement{قضیہ}{#1}}{\par}
\newenvironment{thm}[1][]{\pnbstatement{مسئلہ}{#1}}{\par}
\newenvironment{prob}{\pnbstatement{مشق}{}}{\par}
\newenvironment{proof}{\par\Needspace{4\baselineskip}\textbf{ثبوت۔}\par\nobreak}{\hfill\(\square\)\par}
\newcommand{\pnbref}[2]{#1~\textenglish{\ref{#2}}}
\renewcommand{\chaptername}{باب}
\renewcommand{\figurename}{شکل}
\renewcommand{\contentsname}{فہرست}
\makeatletter
\def\@seccntformat#1{\textenglish{\csname the#1\endcsname}\quad}
\renewcommand{\fnum@figure}{\figurename~\textenglish{\thefigure}}
\makeatother
\begin{document}
\begin{titlepage}
\begin{Arabic}
\centering
\vspace*{18mm}
{\Huge\bfseries سیٹ\par}
\vspace{8mm}
{\LARGE اوپن لاجک\par}
\vspace{5mm}
{\Large پنجابی شاہ مکھی ترجمہ\par}
\vspace{16mm}
ایس جلد وچ سیٹاں دا مکمل باب اے۔ ایہ پوری کتاب دا ترجمہ نہیں۔
\par
ایہ مشینی ترجمہ اے۔ خاص ریاضیاتی اصطلاحاں وچوں کجھ عارضی نیں؛
اوہناں دے ماخذ تے فیصلے نال دتے گئے ریکارڈ وچ ویکھے جا سکدے نیں۔
\vfill
\end{Arabic}
\begin{center}
Open Logic Project\par
Source revision \texttt{9620cc73f9c8e0ad003c514a5d3748f29611c4c0}\par
Punjabi Shahmukhi adaptation by Codex. Independent machine translation.\par
Source coverage: OLP-0004--OLP-0010 (7 of 722 source units).\par
Typeface profile: naskh.\par
\href{https://openlogicproject.org/}{openlogicproject.org}\par
\href{https://github.com/KokunoYumeto/OpenLogic-translations}{OpenLogic translations hub}\par
CC BY 4.0. Translation and layout modified; no upstream endorsement implied.
\end{center}
\end{titlepage}
\mainmatter
\begin{Arabic}
\chapter{سیٹ}\label{sfr:set::chap}



\section{عنصراں نال تعیّن}\label{sfr:set:bas:sec}

\emph{سیٹ} شےواں دا اک اکٹھ اے، جہنوں {اکو} شے منّیا جاندا اے۔ سیٹ بناؤن والیاں
شےواں نوں اوس سیٹ دے \emph{عنصر} یا \emph{رکن} آکھیا جاندا اے۔ جے $x$ سیٹ~$A$ دا
عنصر ہووے، تاں اسی $x \in A$ لکھدے آں؛ جے نہ ہووے، تاں $x \notin A$
لکھدے آں۔ جس سیٹ دے کوئی عنصر نہ ہون، اوہنوں \emph{خالی} سیٹ آکھدے نیں
تے اوہنوں~«$\emptyset$» نال ظاہر کردے نیں۔

\begin{explain}
ایس گل نال کوئی فرق نہیں پیندا کہ اسی سیٹ نوں کِویں \emph{دسدے} آں، اوہدے
عنصراں نوں کِہڑی \emph{ترتیب} وچ رکھدے آں، یا اوہدے عنصراں نوں
\emph{کِنّی واری} گِندے آں۔ اصل گل صرف ایہ اے کہ اوہدے عنصر کِہڑے نیں۔
اسی ایس گل نوں ہِیٹھلے اصول دی صورت وچ بیان کردے آں۔
\end{explain}

\begin{defn}[عنصراں نال تعیّن]
  جے $A$ تے $B$ سیٹ نیں، تاں $A = B$ اُدوں تے صرف اُدوں ہوندا اے جدوں
  $A$ دا ہر عنصر، $B$ دا وی عنصر ہووے، تے ایس توں اُلٹ وی۔
\end{defn}

عنصراں نال تعیّن دا اصول کجھ علامتاں ورتن دا جواز دیندا اے۔ عام طور تے، جدوں
ساڈے کول کجھ شےواں $a_{1}$، \dots، $a_{n}$ ہون، تاں $\{a_{1}, \dots, a_{n}\}$
\emph{اوہی} سیٹ اے جہدے عنصر، $a_1, \ldots, a_n$ نیں۔ اسی لفظ
«\emph{اوہی}» اُتے زور دیندے آں، کیوں جو عنصراں نال تعیّن دا اصول دسدا اے کہ
ایہو جہا سیٹ صرف \emph{اک} ہو سکدا اے۔ ایس اصول توں ایہ وی جائز اے:
  \[
    \{a, a, b\} = \{a, b\} = \{b,a\}.
  \]
ایس نال اوہ گل واضح ہوندی اے کہ سیٹاں بارے سوچدیاں سانوں ایہدی پرواہ نہیں
ہوندی کہ اوہناں دے عنصر دی ترتیب کی اے، یا اوہ کِنّی واری دسے گئے نیں۔

\begin{tagblock}{novice}
\begin{ex}
جدوں وی تہاڈے کول کجھ شےواں ہون، تسیں اوہناں نوں اک سیٹ وچ اکٹھیاں کر سکدے او۔
مثال لئی، رچرڈ دے بھین بھراواں دا سیٹ اک ایہو جہا سیٹ اے جہدے وچ اک بندہ اے،
تے اسی اوہنوں $S=\{\text{\textenglish{Ruth}}\}$ لکھ سکدے آں۔ $4$ توں گھٹ مثبت صحیح
اعداد دا سیٹ $\{1, 2, 3\}$ اے، پر اوہنوں $\{3, 2, 1\}$ یا
$\{1, 2, 1, 2, 3\}$ وی لکھیا جا سکدا اے۔ عنصراں نال تعیّن دے اصول موجب ایہ
سارے اکو سیٹ نیں۔ کیوں جو $\{1, 2, 3\}$ دا ہر عنصر،
$\{3, 2, 1\}$ دا (تے $\{1, 2, 1, 2, 3\}$ دا وی) عنصر اے، تے ایس
توں اُلٹ وی۔
\end{ex}
\end{tagblock}

اکثر اسی سیٹ نوں اوہدے عنصر دی سانجھی خاصیت نال دسّاں گے۔ ایس لئی
اسی ایہ مختصر علامتی لکھت ورتاں گے: $\Setabs{x}{\phi(x)}$، جتھے $\phi(x)$
اوہ خاصیت اے جہڑی $x$ وچ ہونی چاہیدی اے تاں جو اوہ سیٹ دے عنصراں وچوں
اک گنیا جا سکے۔

\begin{tagblock}{novice}
\begin{ex}
اپنی مثال وچ اسی $S$ نوں ایویں وی دس سکدے ساں:
\[
S = \Setabs{x}{\text{\textarabic{\textenglish{\(x\)} رچرڈ دی بھین یا اوہدا بھرا اے}}}.
\]
\end{ex}
\end{tagblock}

\begin{tagblock}{math}
\begin{ex}
کِسے عدد نوں \emph{کامل} اُدوں تے صرف اُدوں آکھدے نیں جدوں اوہ اپنے اوہناں
ونڈن والے اعداد دے جوڑ دے برابر ہووے جو عدد نوں پورا ونڈدے ہون پر آپ اوہدے
برابر نہ ہون۔ مثال لئی، $6$ کامل اے، کیوں جو اوہنوں ونڈن والے اعداد، جو آپ
اوہدے برابر نہیں، $1$، $2$ تے~$3$ نیں، تے $6 = 1 + 2 + 3$ اے۔ دراصل،
$6$، $10$ توں گھٹ اکو اک مثبت صحیح عدد اے جو کامل اے۔ سو، عنصراں نال تعیّن
دا اصول ورت کے اسی آکھ سکدے آں:
\[
	\{6\} = \Setabs{x}{\text{\textarabic{\textenglish{\(x\)} کامل اے تے \textenglish{\(0 \leq x \leq 10\)}}}}
\]
سجے پاسے دی علامتی لکھت نوں اسی ایویں پڑھدے آں: «اوہناں $x$ دا سیٹ جیہناں لئی
$x$ کامل اے تے $0 \leq x \leq 10$»۔ ایتھے برابری دی ایہ لکھت تصدیق کردی
اے کہ سیٹاں بارے سوچدیاں سانوں ایہدی پرواہ نہیں ہوندی کہ اوہناں نوں کِویں
دسیا گیا اے۔ ہور عام طور تے، عنصراں نال تعیّن دا اصول ایہ یقینی بناندا اے کہ
اوہناں $x$ دا سیٹ جنہاں لئی $\phi(x)$ ہووے، ہمیشہ اکو ہوندا اے۔
سو، عنصراں نال تعیّن دا اصول $\Setabs{x}{\phi(x)}$ نوں اوہناں $x$ دا
\emph{اوہی} سیٹ آکھن دا جواز دیندا اے جنہاں لئی~$\phi(x)$ ہووے۔
\end{ex}
\end{tagblock}

عنصراں نال تعیّن دا اصول ایہ وکھاؤن دا طریقہ دیندا اے کہ سیٹ اکو نیں: $A = B$
وکھاؤن لئی ایہ وکھاؤ کہ جدوں وی $x \in A$ ہووے، تاں $x \in B$ وی ہووے،
تے جدوں وی $y \in B$ ہووے، تاں $y \in A$ وی ہووے۔

\begin{prob}
ثابت کرو کہ ودھ توں ودھ اک خالی سیٹ ہو سکدا اے، یعنی ایہ وکھاؤ کہ جے $A$ تے
$B$ ایہو جہے سیٹ نیں جنہاں دے کوئی عنصر نہیں، تاں $A = B$ اے۔
\end{prob}






\section{ذیلی سیٹ تے ذیلی سیٹاں دے سیٹ}\label{sfr:set:sub:sec}

\begin{explain}
سانوں اکثر سیٹاں دا آپس وچ ٹاکرا کرن دی لوڑ پئے گی۔ ٹاکرے دی اک صاف صورت
ایہ اے: \emph{اک سیٹ وچ جیہڑی وی شے اے، اوہ دوجے وچ وی اے}۔ ایہ صورت اینی
اہم اے کہ ایس لئی کجھ نویاں علامتاں متعارف کرائیاں جان۔
\end{explain}

\begin{defn}[ذیلی سیٹ]
جے سیٹ $A$ دا ہر عنصر، $B$ دا وی عنصر ہووے، تاں اسی آکھدے آں
کہ $A$، $B$ دا \emph{ذیلی سیٹ} اے، تے $A \subseteq B$ لکھدے آں۔ جے
$A$، $B$ دا ذیلی سیٹ نہ ہووے، تاں اسی $A \not\subseteq B$ لکھدے آں۔
جے $A \subseteq B$ ہووے پر $A \neq B$ ہووے، تاں اسی $A \subsetneq B$
لکھدے آں تے آکھدے آں کہ $A$، $B$ دا \emph{حقیقی ذیلی سیٹ} اے۔
\end{defn}

\begin{ex}
ہر سیٹ اپنے آپ دا ذیلی سیٹ اے، تے $\emptyset$ ہر سیٹ دا ذیلی سیٹ اے۔ قدرتی
جفت اعداد دا سیٹ قدرتی اعداد دے سیٹ دا ذیلی سیٹ اے۔ ایس طرح
$\{ a, b \} \subseteq \{ a, b, c \}$۔ پر $\{ a, b, e \}$،
$\{ a, b, c \}$ دا ذیلی سیٹ نہیں اے۔
\end{ex}

\begin{ex}
عدد $2$ صحیح اعداد دے سیٹ دا عنصر اے، جد کہ جفت اعداد دا سیٹ صحیح
اعداد دے سیٹ دا ذیلی سیٹ اے۔ پر ایہ وی ہو سکدا اے کہ اک سیٹ کِسے ہور سیٹ
دا عنصر وی ہووے تے اوہدا ذیلی سیٹ وی، یعنی \emph{دونوں} گلاں ہون؛
مثال لئی، $\{0\} \in \{0, \{0\}\}$ تے نال ہی
$\{0\} \subseteq \{0, \{0\}\}$۔
\end{ex}

عنصراں نال تعیّن دا اصول سیٹاں دے اکو ہون دی کسوٹی دیندا اے: $A = B$ اُدوں
تے صرف اُدوں ہوندا اے جدوں $A$ دا ہر عنصر، $B$ دا وی عنصر
ہووے تے ایس توں اُلٹ وی۔ «ذیلی سیٹ» دی تعریف $A \subseteq B$ نوں عین ایس
کسوٹی دے پہلے ادھ دے طور تے دسدی اے: $A$ دا ہر عنصر، $B$ دا وی
عنصر اے۔ بے شک، جے اسی $A$ تے $B$ نوں آپس وچ وٹا دیئیے، تاں وی
تعریف لاگو ہوندی اے: یعنی $B \subseteq A$ اُدوں تے صرف اُدوں جدوں $B$ دا
ہر عنصر، $A$ دا وی عنصر ہووے۔ تے ایہ عین عنصراں نال تعیّن
دے اصول دا «ایس توں اُلٹ وی» والا حصہ اے۔ ہور لفظاں وچ، ایس اصول توں
نکلدا اے کہ سیٹ اُدوں تے صرف اُدوں برابر نیں جدوں اوہ اک دوجے دے ذیلی سیٹ ہون۔

\begin{prop}
$A = B$ اُدوں تے صرف اُدوں ہوندا اے جدوں $A \subseteq B$ تے
$B \subseteq A$، دونوں ہون۔
\end{prop}

ہُن کجھ ہور مددگار علامتاں متعارف کراؤن دا وی چنگا موقع اے۔ ایہ تعریف
کردیاں کہ $A$، $B$ دا ذیلی سیٹ کدوں اے، اسی آکھیا سی کہ «$A$ دا ہر
عنصر \dots اے»، تے «$\dots$» دی تھاں «$B$ دا عنصر»
رکھیا سی۔ پر عبارت دی ایہ \emph{بنت} اینی عام اے کہ ایس لئی کجھ رسمی
علامتاں متعارف کرانا مددگار ہووے گا۔

\begin{defn}\label{sfr:set:sub:forallxina}
$(\forall x \in A)\phi$، $\forall x(x \in A \lif
\phi)$ دی مختصر لکھت اے۔ ایس طرح، $(\exists x \in A)\phi$،
$\exists x(x \in A \land \phi)$ دی مختصر لکھت اے۔
\end{defn}

ایہ علامتاں ورت کے، اسی آکھ سکدے آں کہ $A \subseteq B$ اُدوں تے صرف اُدوں
ہوندا اے جدوں $(\forall x \in A)x \in B$ ہووے۔

ہُن اسی سیٹ دی اک خاص قسم ول آندے آں: کِسے دتے ہوئے سیٹ دے سارے ذیلی
سیٹاں دا سیٹ۔

\begin{defn}[ذیلی سیٹاں دا سیٹ]
سیٹ~$A$ دے سارے ذیلی سیٹاں توں بنے سیٹ نوں $A$ دا \emph{ذیلی سیٹاں دا سیٹ}
آکھدے نیں تے $\Pow{A}$ لکھدے نیں۔
  \[
    \Pow{A} = \Setabs{B}{B \subseteq A}
  \]
\end{defn}

\begin{ex}
$\{ a, b, c \}$ دے سارے ممکن ذیلی سیٹ کِہڑے نیں؟ اوہ ایہ نیں:
$\emptyset$، $\{a \}$، $\{b\}$، $\{c\}$، $\{a, b\}$، $\{a, c\}$،
$\{b, c\}$، $\{a, b, c\}$۔ ایہناں سارے ذیلی سیٹاں دا سیٹ
$\Pow{\{a,b,c\}}$ اے:
\[
\Pow{\{ a, b, c \}} = \{\emptyset, \{a \}, \{b\}, \{c\}, \{a, b\},
\{b, c\}, \{a, c\}, \{a, b, c\}\}
\]
\end{ex}

\begin{prob}
$\{a, b, c, d\}$ دے سارے ذیلی سیٹ لکھو۔
\end{prob}

\begin{prob}
وکھاؤ کہ جے $A$ دے $n$ عنصر ہون، تاں $\Pow{A}$ دے $2^n$
عنصر ہون گے۔
\end{prob}






\section{کجھ اہم سیٹ}\label{sfr:set:imp:sec}

\begin{ex}
ساڈا واسطہ زیادہ تر اوہناں سیٹاں نال پئے گا جیہناں دے عنصر ریاضی
دیاں شےواں نیں۔ ایہو جہے چار سیٹ اینے اہم نیں کہ اوہناں دے خاص ناں نیں:
\begin{multline*}
    \Nat = \{0, 1, 2, 3, \ldots\} \\
    \shoveright{\text{\textarabic{قدرتی اعداد دا سیٹ}}}\\
    \shoveleft{\Int = \{\ldots, -2, -1, 0, 1, 2, \ldots\}} \\
    \shoveright{\text{\textarabic{صحیح اعداد دا سیٹ}}}\\
    \shoveleft{\Rat = \Setabs{\nicefrac{m}{n}}{\text{\textarabic{\textenglish{\(m, n \in \Int\)} تے \textenglish{\(n \neq 0\)}}}}}\\
    \shoveright{\text{\textarabic{ناطق اعداد دا سیٹ}}}\\
    \shoveleft{\Real = (-\infty, \infty)}\\
    \text{\textarabic{حقیقی اعداد دا سیٹ (تسلسل)}}
\end{multline*}
ایہ سارے \emph{لامتناہی} سیٹ نیں، یعنی ہر اک دے عنصر دی گنتی
لامتناہی اے۔

ایہناں سیٹاں وچ اگے ودھدیاں، اسی اپنے کول \emph{ہور} عدد شامل کردے آں۔ ایہ
صاف ہونا چاہیدا اے کہ $\Nat \subseteq \Int \subseteq \Rat \subseteq \Real$:
آخر ہر قدرتی عدد صحیح عدد اے؛ ہر صحیح عدد ناطق عدد اے؛ تے ہر ناطق عدد
حقیقی عدد اے۔ ایس طرح ایہ وی صاف ہونا چاہیدا اے کہ
$\Nat \subsetneq \Int \subsetneq \Rat$، کیوں جو $-1$ صحیح عدد اے پر
قدرتی عدد نہیں، تے $\nicefrac{1}{2}$ ناطق عدد اے پر صحیح عدد نہیں۔ ایہ گل
اینی صاف نہیں کہ $\Rat \subsetneq \Real$، یعنی کجھ حقیقی عدد ایہو جہے
نیں جو ناطق نہیں۔

کدی کدی اسی مثبت صحیح اعداد دا سیٹ $\PosInt = \{1, 2, 3, \dots\}$ تے
صرف پہلے دو قدرتی اعداد والا سیٹ $\Bin = \{0, 1\}$ وی ورتاں گے۔
\end{ex}

\begin{tagblock}{compsci}
\begin{ex}[لڑیاں]
اک ہور دلچسپ مثال سیٹ $A^{*}$ اے، جہڑا حرفاں دے مجموعے $A$ اُتے
\emph{محدود لڑیاں} دا سیٹ اے: $A$ دے عنصراں دا کوئی وی محدود سلسلہ، $A$ اُتے
اک لڑی اے۔ اسی \emph{خالی لڑی $\Lambda$} نوں وی $A$ اُتے بنن والیاں لڑیاں وچ
شامل کردے آں، ہر حرفی مجموعے~$A$ لئی۔ مثال لئی،
\begin{multline*}
\Bin^*
=\{\Lambda,0,1,00,01,10,11,\\
000,001,010,011,100,101,110,111,0000,\ldots\}.
\end{multline*}
جے $x=x_{1}\ldots x_{n}\in A^{*}$ اک لڑی ہووے جہدے وچ $n$ «حرف» ہون،
جو $A$ توں لئے گئے ہون، تاں اسی آکھدے آں کہ لڑی دی \emph{لمبائی}~$n$ اے تے
$\len{x}=n$ لکھدے آں۔
\end{ex}
\end{tagblock}

\begin{ex}[لامتناہی سلسلے]
کِسے وی سیٹ $A$ لئی اسی اوہدے عنصراں دے لامتناہی سلسلیاں
دا سیٹ~$A^\omega$ وی ویکھ سکدے آں، جتھے عنصر $A$ دے نیں۔ اک لامتناہی
سلسلہ $a_1a_2a_3a_4\dots$ شےواں دی اک پاسے نوں لامتناہی فہرست اے، جہدی ہر
شے $A$ دا عنصر اے۔
\end{ex}






\section{اتحاد تے اشتراک}\label{sfr:set:uni:sec}

\begin{explain}
\pnbref{حصہ}{sfr:set:bas:sec} وچ اسی خاصیت راہیں سیٹاں دی تعریف متعارف کرائی
سی، یعنی $\Setabs{x}{\phi(x)}$ دی صورت والیاں تعریفاں۔ ایتھے اسی کوئی
خاصیت~$\phi$ ورتدے آں، تے ایس خاصیت وچ اوہناں سیٹاں دا ذکر ہو سکدا اے
جنہاں دی تعریف اسی پہلاں کر چکے آں۔ مثال لئی، جے $A$ تے~$B$ سیٹ نیں،
تاں سیٹ $\Setabs{x}{x \in A \lor x \in B}$ اوہناں ساریاں شےواں دا بنیا اے
جو $A$ یا~$B$ دے عنصر نیں؛ یعنی ایہ اوہ سیٹ اے جو $A$ تے~$B$ دے
عنصراں نوں اکٹھے کردا اے۔ اسی ایہنوں \pnbref{شکل}{sfr:set:uni:fig:union}
وانگ شکل وچ ویکھ سکدے آں، جتھے اُبھاریا ہویا علاقہ دوہاں سیٹاں $A$ تے~$B$
دے عنصراں نوں اکٹھے ظاہر کردا اے۔

\begin{figure}
  \begin{LTR}\centering\input{C:/interlanguage-production/openlogic-pnb-Arab-PK/repo/upstream/assets/diagrams/union.tikz}\end{LTR}
  \caption{دو سیٹاں دا اتحاد $A \cup B$، $A$ تے~$B$ دے
   عنصراں نوں اکٹھے لے کے بنن والا سیٹ اے۔}
  \label{sfr:set:uni:fig:union}
\end{figure}

سیٹاں اُتے ایہ عمل---اوہناں نوں اکٹھا کرنا---بہت فائدےمند تے عام اے، ایس
لئی اسی ایہنوں اک رسمی ناں تے اک علامت دیندے آں۔
\end{explain}

\begin{defn}[اتحاد]
دو سیٹاں $A$ تے $B$ دا \emph{اتحاد}، جہنوں $A \cup B$ لکھدے نیں،
اوہناں ساریاں شےواں دا سیٹ اے جو $A$، $B$ یا دوہاں دے عنصر نیں۔
\[
A \cup B = \Setabs{x}{x \in A \lor x \in B}
\]
\end{defn}

\begin{ex}
کیوں جو عنصراں نوں وار وار گِنن دا کوئی فرق نہیں پیندا،
دو ایہو جہے سیٹاں دا اتحاد جنہاں دا عنصر سانجھا ہووے، اوس
عنصر نوں صرف اک واری رکھدا اے؛ مثال لئی،
$\{ a, b, c\} \cup \{ a, 0, 1\} = \{a, b, c, 0, 1\}$۔

کِسے سیٹ تے اوہدے کِسے ذیلی سیٹ دا اتحاد بس وڈا سیٹ ای اے:
$\{a, b, c \} \cup \{a \} = \{a, b, c\}$۔

سیٹ دا خالی سیٹ نال اتحاد اوہی سیٹ اے:
$\{a, b, c \} \cup \emptyset = \{a, b, c \}$۔
\end{ex}

\begin{prob}
ثابت کرو کہ جے $A \subseteq B$ ہووے، تاں $A \cup B = B$ اے۔
\end{prob}

\begin{explain}
اسی اتحاد دے «دوہرے ہم صورت» عمل اُتے وی غور کر سکدے آں۔ ایہ اوہ عمل
اے جو سارے اوہناں عنصراں دا سیٹ بناندا اے جو
$A$ دے عنصر وی ہون تے $B$ دے عنصر وی۔ ایس عمل نوں
\emph{اشتراک} آکھدے نیں، تے ایہنوں \pnbref{شکل}{sfr:set:uni:fig:intersection} وانگ وکھایا
جا سکدا اے۔
\begin{figure}
  \begin{LTR}\centering\input{C:/interlanguage-production/openlogic-pnb-Arab-PK/repo/upstream/assets/diagrams/intersection.tikz}\end{LTR}
  \caption{دو سیٹاں دا اشتراک $A \cap B$ اوہناں دے سانجھے
    عنصراں دا سیٹ اے۔}
  \label{sfr:set:uni:fig:intersection}
\end{figure}
\end{explain}

\begin{defn}[اشتراک]
دو سیٹاں $A$ تے $B$ دا \emph{اشتراک}، جہنوں $A \cap B$ لکھدے نیں،
اوہناں ساریاں شےواں دا سیٹ اے جو $A$ تے~$B$ دوہاں دے عنصر نیں۔
\[
A \cap B = \Setabs{x}{x \in A \land x \in B}
\]
دو سیٹاں نوں \emph{بے سانجھ} آکھدے نیں جے اوہناں دا اشتراک خالی ہووے۔
ایس دا مطلب اے کہ اوہناں دے کوئی عنصر سانجھے نہیں۔
\end{defn}

\begin{ex}
جے دو سیٹاں دے کوئی عنصر سانجھے نہ ہون، تاں اوہناں دا اشتراک
خالی اے: $\{ a, b, c\} \cap \{ 0, 1\} = \emptyset$۔

جے دو سیٹاں دے کجھ عنصر سانجھے ہون، تاں اوہناں دا اشتراک اوہناں
ساریاں دا سیٹ اے: \[
\{a, b, c \} \cap \{a, b, d \} = \{a, b\}\text{\textarabic{۔}}
\]

کِسے سیٹ دا اوہدے کِسے ذیلی سیٹ نال اشتراک بس چھوٹا سیٹ ای اے:
$\{a, b, c\} \cap \{a, b\} = \{a, b\}$۔

کِسے وی سیٹ دا خالی سیٹ نال اشتراک خالی اے:
$\{a, b, c \} \cap \emptyset = \emptyset$۔
\end{ex}

\begin{prob}
پوری منطقی پابندی نال ثابت کرو کہ جے $A \subseteq B$ ہووے، تاں
$A \cap B = A$ اے۔
\end{prob}

\begin{explain}
اسی دو توں ودھ سیٹاں دا اتحاد یا اشتراک وی بنا سکدے آں۔ عام صورت وچ ایس
لئی اک سوہنا طریقہ ایہ اے: منّ لو تسیں اوہناں سارے سیٹاں نوں اکو سیٹ
وچ اکٹھا کر لیندے او جنہاں دا اتحاد (یا اشتراک) بنانا اے۔ فیر اسی اپنے
اصل سیٹاں دے اتحاد نوں اوہناں ساریاں شےواں دے سیٹ دے طور تے متعین کر
سکدے آں جو اکٹھے کیتے سیٹ دے گھٹ توں گھٹ اک عنصر وچ ہون، تے
اشتراک نوں اوہناں ساریاں شےواں دے سیٹ دے طور تے جو اوہدے ہر عنصر
وچ ہون۔
\end{explain}

\begin{defn}
جے $A$ سیٹاں دا سیٹ اے، تاں $\bigcup A$، سیٹ~$A$ دے
عنصراں دے عنصراں دا سیٹ اے:
\begin{align*}
\bigcup A & = \Setabs{x}{\text{\textarabic{\textenglish{\(x\)}، سیٹ \textenglish{\(A\)} دے کسے عنصر دا عنصر اے}}},
\text{\textarabic{ یعنی،}}\\
& = \Setabs{x}{\text{\textarabic{کوئی \textenglish{\(B \in A\)} ایہو جہا اے کہ \textenglish{\(x \in B\)}}}}
\end{align*}
\end{defn}

\begin{defn}
جے $A$ سیٹاں دا سیٹ اے، تاں $\bigcap A$ اوہناں شےواں دا سیٹ اے جو
$A$ دے سارے عنصراں وچ سانجھیاں نیں:
\begin{align*}
\bigcap A & = \Setabs{x}{\text{\textarabic{\textenglish{\(x\)}، سیٹ \textenglish{\(A\)} دے ہر عنصر وچ اے}}},
\text{\textarabic{ یعنی،}}\\
 & = \Setabs{x}{\text{\textarabic{ہر \textenglish{\(B \in A, x \in B\)}}}}
\end{align*}
\end{defn}

\begin{ex}
منّ لو کہ $A = \{ \{ a, b \}, \{ a, d, e \}, \{ a, d \} \}$۔
فیر $\bigcup A = \{ a, b, d, e \}$ تے $\bigcap A = \{ a \}$ اے۔
\end{ex}
\begin{prob}
	وکھاؤ کہ جے $A$ اک سیٹ اے تے $A \in B$، تاں $A \subseteq \bigcup B$ اے۔
\end{prob}

اسی سیٹاں دے سلسلے $A_1$، $A_2$، \dots لئی وی ایہو کم کر سکدے آں:
\begin{align*}
\bigcup_i A_i & = \Setabs{x}{\text{\textarabic{\textenglish{\(x\)}، سیٹاں \textenglish{\(A_i\)} وچوں کسے اک وچ شامل اے}}}\\
\bigcap_i A_i & = \Setabs{x}{\text{\textarabic{\textenglish{\(x\)}، سیٹاں \textenglish{\(A_i\)} وچوں ہر اک وچ شامل اے}}}.
\end{align*}

جدوں ساڈے کول سیٹاں دا \emph{اشاریہ} ہووے، یعنی کوئی سیٹ $I$ جس لئی
اسی $A_i$ اُتے غور کردے ہاں، ہر $i \in I$ لئی، تاں اسی ایہ مختصر
لکھتاں وی ورت سکدے آں:
\begin{align*}
	\bigcup_{i \in I} A_i & = \bigcup \Setabs{A_i }{i \in I}\\
	\bigcap_{i \in I} A_i & = \bigcap\Setabs{A_i}{i \in I}
\end{align*}

آخر وچ، اسی سارے اوہناں عنصراں دے سیٹ بارے سوچ سکدے
آں جو $A$ وچ نیں پر $B$ وچ نہیں۔ ایہنوں \pnbref{شکل}{sfr:set:uni:difference} وانگ وکھایا
جا سکدا اے۔

\begin{figure}
  \begin{LTR}\centering\input{C:/interlanguage-production/openlogic-pnb-Arab-PK/repo/upstream/assets/diagrams/difference.tikz}\end{LTR}
  \caption{دو سیٹاں دا فرق $A \setminus B$ اوہناں عنصراں
    دا سیٹ اے جو $A$ وچ نیں پر $B$ دے عنصر نہیں۔}
  \label{sfr:set:uni:difference}
\end{figure}

\begin{defn}[فرق]
\emph{سیٹاں دا فرق}~$A \setminus B$ سارے اوہناں عنصراں
دا سیٹ اے جو $A$ وچ نیں پر $B$ دے عنصر نہیں، یعنی
\[
A\setminus B = \Setabs{x}{\text{\textarabic{\textenglish{\(x\in A\)} تے \textenglish{\(x \notin B\)}}}}.
\]
\end{defn}

\begin{prob}
	ثابت کرو کہ جے $A \subsetneq B$ ہووے، تاں $B \setminus A \neq \emptyset$ اے۔
\end{prob}






\section{جوڑے، ٹپل تے کارتیسی حاصلِ ضرب}\label{sfr:set:pai:sec}

\begin{explain}
عنصراں نال تعیّن دے اصول توں نکلدا اے کہ سیٹ دے عنصراں دی کوئی ترتیب نہیں
ہوندی۔ سو، جے اسی ترتیب ظاہر کرنا چاہیئے، تاں \emph{ترتیب وار جوڑے}
$\tuple{x, y}$ ورتدے آں۔ بے ترتیب جوڑے $\{x, y\}$ وچ ترتیب دا کوئی فرق
نہیں پیندا: $\{x, y\} = \{y, x\}$۔ ترتیب وار جوڑے وچ فرق پیندا اے: جے
$x \neq y$ ہووے، تاں $\tuple{x, y} \neq \tuple{y, x}$ اے۔

سیٹاں دے نظریے وچ ترتیب وار جوڑیاں بارے ساڈا تصور کی ہونا چاہیدا اے؟ بنیادی
گل ایہ اے کہ اسی ایہ خیال قائم رکھنا چاہندے آں کہ ترتیب وار جوڑے اُدوں تے
صرف اُدوں اکو ہوندے نیں جدوں اوہناں دا پہلا عنصر اکو ہووے تے دوجا عنصر وی
اکو ہووے، یعنی:
\[
  \tuple{a, b}= \tuple{c, d}\text{\textarabic{ اُدوں تے صرف اُدوں، جدوں دونوں }}a = c \text{\textarabic{ تے }}b=d.
\]
سیٹاں دے نظریے وچ اسی وینر--کوراتوسکی دی تعریف نال ترتیب وار جوڑے دی
تعریف کر سکدے آں۔
\end{explain}

\begin{defn}[ترتیب وار جوڑا]\label{sfr:set:pai:wienerkuratowski}
	$\tuple{a, b} = \{\{a\}, \{a, b\}\}$۔
\end{defn}

\begin{prob}
	\pnbref{تعریف}{sfr:set:pai:wienerkuratowski} ورت کے ثابت کرو کہ
	$\tuple{a, b}= \tuple{c, d}$ اُدوں تے صرف اُدوں ہوندا اے جدوں
	$a = c$ تے $b=d$، دونوں ہون۔
\end{prob}

\begin{explain}
ترتیب وار جوڑے دی تعریف طے کر کے اسی اوہنوں ہور سیٹاں دی تعریف لئی ورت
سکدے آں۔ مثال لئی، کدی کدی سانوں دو توں ودھ شےواں دے ترتیب وار سلسلے وی
چاہیدے ہوندے نیں، جیویں \emph{تِکے} $\tuple{x, y, z}$، \emph{چوکے}
$\tuple{x, y, z, u}$، تے ایس طرح اگے۔ تِکیاں نوں اسی خاص ترتیب وار جوڑے
سمجھ سکدے آں، جنہاں دا پہلا عنصر آپ اک ترتیب وار جوڑا اے:
$\tuple{x, y, z}$، $\tuple{\tuple{x, y},z}$ اے۔ چوکیاں لئی وی ایہو گل اے:
$\tuple{x,y,z,u}$، $\tuple{\tuple{\tuple{x,y},z},u}$ اے، تے ایس طرح اگے۔
عام طور تے اسی \emph{ترتیب وار $n$-ٹپل} $\tuple{x_1, \dots, x_n}$ دی گل
کردے آں۔

ترتیب وار جوڑیاں، یا ہور ترتیب وار $n$-ٹپلاں دے کجھ سیٹ ساڈے کم آؤن گے۔
\end{explain}

\begin{defn}[کارتیسی حاصلِ ضرب]
دتے ہوئے سیٹاں $A$ تے $B$ دا \emph{کارتیسی حاصلِ ضرب} $A \times B$ ایویں
متعین ہوندا اے:
\[
  A \times B = \Setabs{\tuple{x, y}}{\text{\textarabic{\textenglish{\(x \in A\)} تے \textenglish{\(y \in B\)}}}}.
\]
\end{defn}

\begin{ex}
جے $A = \{0, 1\}$ تے $B = \{1, a, b\}$ ہون، تاں اوہناں دا حاصلِ ضرب اے:
\[
A \times B = \{ \tuple{0, 1}, \tuple{0, a}, \tuple{0, b},
    \tuple{1, 1}, \tuple{1, a}, \tuple{1, b} \}.
\]
\end{ex}

\begin{ex}
جے $A$ اک سیٹ اے، تاں $A$ دا اپنے آپ نال حاصلِ ضرب، $A \times A$،
$A^2$ وی لکھیا جاندا اے۔ ایہ اوہناں \emph{سارے} جوڑیاں $\tuple{x, y}$ دا
سیٹ اے جنہاں لئی $x, y \in A$ اے۔ سارے تِکیاں $\tuple{x, y, z}$ دا
سیٹ $A^3$ اے، تے ایس طرح اگے۔ اسی اک بازگشتی تعریف دے سکدے آں:
\begin{align*}
  A^1 & = A\\
  A^{k+1} & = A^k \times A
\end{align*}
\end{ex}

\begin{prob}
سارے عنصر لکھو جو $\{1, 2, 3\}^3$ وچ نیں۔
\end{prob}

\begin{prop}\label{sfr:set:pai:cardnmprod}
جے $A$ دے $n$ عنصر تے $B$ دے $m$ عنصر ہون، تاں
$A \times B$ دے $n\cdot m$ عنصر ہون گے۔
\end{prop}

\begin{proof}
ہر عنصر~$x$ لئی جو $A$ وچ اے، $m$ ایہو جہے عنصر نیں
جنہاں دی صورت $\tuple{x, y} \in A \times B$ اے۔ منّ لو کہ
$B_x = \Setabs{\tuple{x, y}}{y \in B}$۔ کیوں جو جدوں وی $x_1 \neq x_2$
ہووے، تاں $\tuple{x_1, y} \neq \tuple{x_2, y}$ اے، ایس لئی
$B_{x_1} \cap B_{x_2} = \emptyset$۔ پر جے $A = \{x_1, \dots, x_n\}$
ہووے، تاں $A \times B = B_{x_1} \cup \dots \cup B_{x_n}$ اے، تے ایس لئی
اوہدے $n\cdot m$ عنصر نیں۔

ایس گل نوں شکل وچ ویکھن لئی، عنصراں نوں جو $A \times B$
وچ نیں، قطاراں تے کالماں دے جال وچ ترتیب نال رکھو:
\[
\begin{array}{rcccc}
  B_{x_1} = & \{\tuple{x_1, y_1} & \tuple{x_1, y_2} & \dots & \tuple{x_1, y_m}\}\\
  B_{x_2} = & \{\tuple{x_2, y_1} & \tuple{x_2, y_2} & \dots & \tuple{x_2, y_m}\}\\
  \vdots & & \vdots\\
  B_{x_n} = & \{\tuple{x_n, y_1} & \tuple{x_n, y_2} & \dots & \tuple{x_n, y_m}\}
\end{array}
\]
کیوں جو سارے $x_i$ وکھ وکھ نیں، تے سارے $y_j$ وی وکھ وکھ نیں، ایس جال
وچ کوئی دو جوڑے اکو نہیں، تے اوہناں دی گنتی $n\cdot m$ اے۔
\end{proof}

\begin{prob}
$k$ اُتے ریاضیاتی استقرا نال وکھاؤ کہ سارے $k \ge 1$ لئی، جے $A$ دے
$n$ عنصر ہون، تاں $A^k$ دے $n^k$ عنصر ہون گے۔
\end{prob}

\begin{ex}
جے $A$ اک سیٹ اے، تاں $A$ اُتے \emph{لفظ}، عنصراں دا
کوئی وی سلسلہ اے جو $A$ توں لئے گئے ہون۔ سلسلے نوں اک $n$-ٹپل سمجھیا جا
سکدا اے، جہدے عنصر، $A$ توں نیں۔ مثال لئی، جے $A = \{a, b, c\}$
ہووے، تاں سلسلے «$bac$» نوں تِکّے~$\tuple{b, a, c}$ دے طور تے سمجھیا
جا سکدا اے۔ لفظ، یعنی علامتاں دے سلسلے، کمپیوٹر سائنس وچ بنیادی اہمیت
رکھدے نیں۔ رواج موجب، اسی عنصراں نوں جو $A$ وچ نیں،
لمبائی~$1$ والے سلسلے گِندے آں، تے $\emptyset$ نوں لمبائی~$0$ والا سلسلہ۔
فیر $A$ اُتے \emph{سارے} لفظاں دا سیٹ اے:
\[
A^* = \{\emptyset\} \cup A \cup A^2 \cup A^3 \cup \dots
\]
\end{ex}






\section{رسل دا تضاد}\label{sfr:set:rus:sec}

عنصراں نال تعیّن دا اصول علامتی لکھت $\Setabs{x}{\phi(x)}$ دا جواز دیندا
اے، جہڑی اوہناں $x$ دا \emph{اوہی} سیٹ دسدی اے جنہاں لئی~$\phi(x)$
ہووے۔ پر عنصراں نال تعیّن دا اصول \emph{اصل وچ} صرف ایس خیال دا جواز
دیندا اے: \emph{جے} اک ایہو جہا سیٹ موجود ہووے جہدے رکن سارے اوہ تے صرف
اوہ ہون جیہڑے $\phi$ نیں، \emph{تاں} ایہو جہا سیٹ اکو اے۔ ہور لفظاں وچ:
کِسے~$\phi$ نوں طے کرن مگروں، سیٹ $\Setabs{x}{\phi(x)}$ یکتا اے،
\emph{جے اوہ موجود ہووے}۔

پر ایہ شرط بڑی اہم اے!{} بنیادی گل ایہ اے کہ ہر خاصیت توں
\emph{سیٹ نہیں بنایا جا سکدا}۔ یعنی کجھ خاصیتاں سیٹ دی تعریف
\emph{نہیں} کردیاں۔ جے ساریاں کردیاں، تاں اسی صاف تضاداں وچ پھس جاواں
گے۔ ایس دی سب توں مشہور مثال رسل دا تضاد اے۔

سیٹ ہور سیٹاں دے عنصر ہو سکدے نیں---مثال لئی، سیٹ~$A$ دا ذیلی
سیٹاں دا سیٹ، سیٹاں توں ای بنیا ہوندا اے۔ سو ایہ پچھنا یا کھوجنا بجا اے
کہ کوئی سیٹ دوجے سیٹ دا عنصر اے یا نہیں۔ کی کوئی سیٹ اپنے آپ دا
رکن ہو سکدا اے؟ سیٹ دے خیال وچ کوئی گل ایس نوں روکندی نہیں لگدی۔ مثال
لئی، جے \emph{سارے} سیٹ شےواں دا اک اکٹھ بناندے نیں، تاں کوئی ایہ سوچ
سکدا اے کہ اوہناں نوں اکو سیٹ وچ اکٹھا کیتا جا سکدا اے---سارے سیٹاں دا
سیٹ۔ تے اوہ، آپ اک سیٹ ہون پاروں، سارے سیٹاں دے سیٹ دا عنصر
ہووے گا۔

رسل دا تضاد اُدوں پیدا ہوندا اے جدوں اسی اپنے آپ نوں عنصر نہ
رکھن دی خاصیت، یعنی \emph{اپنے آپ دا رکن نہ ہون} دی خاصیت، اُتے غور کردے
آں۔ جے اسی منّ لئیے کہ اوہناں سارے سیٹاں دا سیٹ موجود اے جو اپنے آپ نوں
عنصر نہیں رکھدے، تاں کی ہووے گا؟ کی
\[
R = \Setabs{x}{x \notin x}
\]
موجود اے؟ گل ایہ اے کہ اسی ثابت کر سکدے آں کہ ایہ موجود نہیں۔

\begin{thm}[رسل دا تضاد]\label{sfr:set:rus:thm:russells-paradox}
	کوئی سیٹ $R = \Setabs{x}{x \notin x}$ موجود نہیں۔
\end{thm}

\begin{proof}
جے $R = \Setabs{x}{x \notin x}$ موجود ہووے، تاں
$R \in R$ اُدوں تے صرف اُدوں جدوں $R \notin R$، تے ایہ اک تضاد اے۔
\end{proof}

\begin{tagblock}{novice}
\begin{explain}
آؤ ایس ثبوت نوں ہولی ہولی ویکھیئے۔ جے $R$ موجود اے، تاں ایہ پچھنا بجا اے
کہ $R \in R$ اے یا نہیں۔ منّ لو کہ واقعی $R \in R$ اے۔ ہُن $R$ دی
تعریف اوہناں سارے سیٹاں دے سیٹ دے طور تے کیتی گئی سی جو اپنے آپ دے
عنصر نہیں۔ سو، جے $R \in R$، تاں $R$ آپ اوہ خاصیت نہیں رکھدا
جس نال $R$ دی تعریف کیتی گئی اے۔ پر صرف اوہ سیٹ $R$ وچ نیں جو ایہ
خاصیت رکھدے نیں، ایس لئی $R$، $R$ دا عنصر نہیں ہو سکدا، یعنی
$R \notin R$۔ پر $R$، $R$ دا عنصر ہووے وی تے نہ وی ہووے، ایہ
نہیں ہو سکدا؛ سو ساڈے کول اک تضاد اے۔

کیوں جو ایہ فرض کرنا کہ $R \in R$ تضاد ول لے جاندا اے، ایس لئی
$R \notin R$ اے۔ پر ایہ وی تضاد ول لے جاندا اے!{} کیوں جو جے $R \notin R$
ہووے، تاں $R$ آپ اوہ خاصیت رکھدا اے جس نال $R$ دی تعریف کیتی گئی اے،
تے ایس لئی $R$، $R$ دا عنصر ہووے گا، بالکل اوہناں ہور سیٹاں
وانگ جو اپنے آپ دے رکن نہیں۔ تے فیر، ایہ نہیں ہو سکدا کہ اوہ $R$ دا
عنصر نہ وی ہووے تے ہووے وی۔
\end{explain}
\end{tagblock}

\begin{digress}
اسی سیٹاں دا ایہو جہا نظریہ کِویں بنائیے جو رسل دے تضاد وچ نہ پھسے، یعنی
ایہ \emph{بے مطابقت} دعویٰ نہ کرے کہ $R = \Setabs{x}{x \notin x}$
موجود اے؟ ایس لئی سانوں مسلّمے مقرر کرنے پَین گے، جو ایہ دسن لئی بالکل
درست شرطاں دین کہ سیٹ کدوں موجود ہوندے نیں (تے کدوں نہیں ہوندے)۔

ایس باب وچ سیٹاں دے نظریے دا جیہڑا خاکہ دتا گیا اے، اوہ ایہ کم نہیں کردا۔
ایہ \emph{واقعی سادہ فہم} نظریہ اے۔ ایہ صرف دسدا اے کہ سیٹ عنصراں نال
تعیّن دے اصول اُتے چلدے نیں تے، جے تہاڈے کول کجھ سیٹ ہون، تاں تسیں اوہناں
دا اتحاد، اشتراک وغیرہ بنا سکدے او۔ سیٹاں دے نظریے نوں ایس توں زیادہ
سخت منطقی بنیاد اُتے اساریا جا سکدا اے۔ 
\end{digress}



\end{Arabic}
\end{document}
